\documentclass[a4paper,12pt]{article}
\usepackage[T1]{fontenc}
\usepackage[utf8]{inputenc}
\usepackage[polish]{babel}
\usepackage{geometry}
\usepackage{array}
\usepackage{longtable}
\usepackage{graphicx}
\usepackage{hyperref}
\geometry{margin=2.5cm}

\title{Sprawozdanie końcowe\\Projekt \textit{Senior Monitor}}
\author{Zespół projektowy}
\date{\today}

\begin{document}
\maketitle
\tableofcontents
\newpage

\section{Wstęp}
Celem projektu \textbf{Senior Monitor} było stworzenie systemu do zdalnego monitorowania podstawowych parametrów życiowych seniora oraz jego lokalizacji. System został zaprojektowany jako rozwiązanie IoT składające się z warstwy urządzenia (ESP32 + czujniki), warstwy komunikacyjnej (MQTT), warstwy serwerowej (Node.js) oraz aplikacji webowej (React/Vite).

Zakładane i zrealizowane funkcjonalności:
\begin{itemize}
  \item cykliczny odczyt danych z wielu czujników medycznych i środowiskowych,
  \item publikacja danych telemetrycznych przez MQTT w czasie zbliżonym do rzeczywistego,
  \item wykrywanie sytuacji alarmowych (np. możliwy upadek, hipoksja, nieprawidłowe tętno, przycisk SOS),
  \item zapisywanie alertów i lokalizacji w bazie danych Supabase,
  \item podgląd danych na frontendzie, mapa pozycji GPS i panel alertów,
  \item zdalne sterowanie urządzeniem (LED, buzzer) z poziomu panelu administracyjnego.
\end{itemize}

\section{Wstęp teoretyczny}
Projekt opiera się na klasycznym modelu \textbf{Internetu Rzeczy (IoT)}:
\begin{enumerate}
  \item \textbf{Warstwa akwizycji danych} --- czujniki zbierają sygnały fizyczne i biologiczne.
  \item \textbf{Warstwa transportowa} --- dane są przesyłane protokołem MQTT.
  \item \textbf{Warstwa przetwarzania i składowania} --- backend filtruje dane, wykrywa alarmy i zapisuje rekordy do bazy.
  \item \textbf{Warstwa prezentacji} --- aplikacja webowa wizualizuje stan pacjenta i urządzenia.
\end{enumerate}

\textbf{MQTT} jest lekkim protokołem publish/subscribe, dobrze dopasowanym do urządzeń o ograniczonych zasobach. W projekcie:
\begin{itemize}
  \item ESP32 publikuje dane na topicach \texttt{STM\_<id>/...},
  \item frontend i backend subskrybują dane przez brokera Mosquitto,
  \item frontend dodatkowo wysyła komendy sterujące na topici \texttt{ESP\_<id>/...}.
\end{itemize}

Do składowania danych wykorzystano \textbf{Supabase} (PostgreSQL + API), co upraszcza integrację z frontendem i backendem.

\section{Dokumentacja urządzenia (ESP32 + czujniki)}
\subsection{Platforma}
Urządzenie bazowe: \textbf{ESP32} z firmwarem MicroPython. Kod główny znajduje się w pliku \texttt{esp32/main.py}.

\subsection{Zastosowane czujniki i elementy}
\begin{itemize}
  \item DS18B20 --- temperatura (1-Wire),
  \item DHT11 --- temperatura i wilgotność,
  \item MPU6050 --- akcelerometr + żyroskop,
  \item SW-520D --- czujnik drgań/nachylenia,
  \item MAX30102 --- tętno i saturacja krwi (SpO\textsubscript{2}),
  \item GY-NEO6MV2 --- moduł GPS,
  \item KY-037 --- czujnik dźwięku (odczyt analogowy),
  \item przycisk SOS,
  \item buzzer,
  \item diody LED: zielona i czerwona.
\end{itemize}

\subsection{Mapowanie pinów (aktualna konfiguracja firmware)}
\begin{longtable}{|>{\raggedright\arraybackslash}p{5cm}|p{3cm}|p{6cm}|}
\hline
\textbf{Funkcja} & \textbf{GPIO} & \textbf{Uwagi} \\
\hline
DS18B20 & 2 & Magistrala 1-Wire \\
\hline
SW-520D & 3 & Odczyt cyfrowy \\
\hline
I2C SDA / SCL & 4 / 5 & MPU6050, MAX30102 \\
\hline
DHT11 & 6 & Odczyt temp./wilgotności \\
\hline
Przycisk SOS & 7 & Tryb pull-up \\
\hline
LED czerwona / zielona & 8 / 10 & Zdalne sterowanie przez MQTT \\
\hline
KY-037 (ADC) & 0 & Odczyt analogowy \\
\hline
Buzzer & 1 & Sterowanie czasem brzęczenia (ms) \\
\hline
GPS UART RX/TX & 20 / 21 & NEO-6M, 9600 bps \\
\hline
\end{longtable}

\subsection{Publikowane dane MQTT}
Każdy odczyt publikowany jest do tematu \texttt{STM\_001/<sensor>}. Przykłady:
\begin{itemize}
  \item \texttt{STM\_001/mpu6050} --- JSON \texttt{\{ax, ay, az, gx, gy, gz, temp\}},
  \item \texttt{STM\_001/dht11} --- JSON \texttt{\{temp, hum\}},
  \item \texttt{STM\_001/max30102\_bpm}, \texttt{STM\_001/max30102\_spo2},
  \item \texttt{STM\_001/neo6m} --- JSON z danymi pozycji GPS,
  \item \texttt{STM\_001/button}, \texttt{STM\_001/sw520d}, \texttt{STM\_001/ky037}.
\end{itemize}

\section{Dokumentacja warstwy serwerowej i bazy danych}
\subsection{Warstwa serwerowa (Node.js + Express)}
Backend znajduje się w \texttt{server/server.js}. Główne zadania serwera:
\begin{itemize}
  \item subskrypcja MQTT (\texttt{+/+/\#}) i odbiór telemetrii ze wszystkich urządzeń,
  \item parsowanie payloadów oraz identyfikacja urządzenia po prefiksie \texttt{STM\_} / \texttt{ESP\_},
  \item detekcja alertów medycznych i zdarzeń krytycznych,
  \item ograniczanie częstości zapisu danych do bazy (interwały czasowe),
  \item udostępnienie API ustawień urządzenia:
  \begin{itemize}
    \item \texttt{GET /api/settings/:deviceId}
    \item \texttt{POST /api/settings/:deviceId}
  \end{itemize}
\end{itemize}

\subsection{Reguły alertowe}
Przykładowe warunki alarmowe zaimplementowane w backendzie:
\begin{itemize}
  \item tętno poniżej 50 lub powyżej 150 ud./min,
  \item saturacja poniżej 90\%,
  \item temperatura ciała poza zakresem bezpiecznym,
  \item naciśnięcie przycisku SOS,
  \item możliwy upadek (analiza sekwencji z MPU6050).
\end{itemize}

\subsection{Baza danych (Supabase/PostgreSQL)}
W projekcie wykorzystywane są co najmniej dwie tabele:
\begin{itemize}
  \item \textbf{alerts} --- rekordy alertów (typ, komunikat, czas, status potwierdzenia),
  \item \textbf{locations} --- historia pozycji GPS (device\_id, latitude, longitude, created\_at).
\end{itemize}

Dodatkowo ustawienia aktywnych czujników dla urządzeń utrzymywane są lokalnie w pliku \texttt{server/device\_settings.json}.

\section{Dokumentacja frontendu wraz z prezentacją}
\subsection{Technologie i struktura}
Frontend znajduje się w katalogu \texttt{pwa/} i jest oparty o:
\begin{itemize}
  \item React + Vite,
  \item \texttt{mqtt.js} (odbiór telemetrii po WebSocket, port 9001),
  \item Supabase JS (pobieranie alertów i lokalizacji),
  \item React Leaflet (wizualizacja mapy).
\end{itemize}

Główne moduły UI:
\begin{itemize}
  \item mapa lokalizacji seniora (LIVE + fallback z historii),
  \item panel czujników (widok admin),
  \item uproszczony panel parametrów życiowych (widok użytkownika),
  \item panel alertów i potwierdzania alarmów,
  \item panel ustawień aktywności czujników,
  \item panel sterowania urządzeniem (LED/buzzer).
\end{itemize}

\subsection{Placeholdery do prezentacji ekranów}
\textbf{Uwaga:} poniższe ramki są miejscem na finalne zrzuty ekranu.

\begin{figure}[h]
\centering
\fbox{\parbox[c][5cm][c]{0.9\linewidth}{\centering Placeholder: ekran główny aplikacji (dashboard)}}
\caption{Widok główny frontendu}
\end{figure}

\begin{figure}[h]
\centering
\fbox{\parbox[c][5cm][c]{0.9\linewidth}{\centering Placeholder: mapa GPS i status sygnału}}
\caption{Moduł mapy lokalizacji}
\end{figure}

\begin{figure}[h]
\centering
\fbox{\parbox[c][5cm][c]{0.9\linewidth}{\centering Placeholder: panel alertów i przycisk potwierdzenia}}
\caption{Moduł alertów}
\end{figure}

\begin{figure}[h]
\centering
\fbox{\parbox[c][5cm][c]{0.9\linewidth}{\centering Placeholder: panel ustawień i sterowania urządzeniem}}
\caption{Ustawienia i sterowanie}
\end{figure}

\section{Podsumowanie}
Projekt \textbf{Senior Monitor} zrealizował główny cel: dostarczenie działającego prototypu systemu monitorowania seniora w czasie (prawie) rzeczywistym. Integracja ESP32, MQTT, backendu Node.js, bazy Supabase i frontendu React umożliwiła stworzenie spójnego rozwiązania obejmującego pomiar, transmisję, analizę i prezentację danych.

Najważniejsze osiągnięcia:
\begin{itemize}
  \item stabilny pipeline telemetrii IoT od czujnika do interfejsu webowego,
  \item mechanizm alertowania i potwierdzania zdarzeń alarmowych,
  \item zapis historii lokalizacji i zdarzeń w bazie danych,
  \item podział uprawnień i trybów pracy frontendu (admin/użytkownik).
\end{itemize}

Rekomendowane kierunki rozwoju:
\begin{itemize}
  \item pełne uwierzytelnianie i autoryzacja dostępu do MQTT oraz API,
  \item szyfrowanie TLS dla komunikacji MQTT/HTTP,
  \item automatyczne raporty medyczne i analiza trendów długoterminowych,
  \item wdrożenie produkcyjne frontendu (build statyczny + reverse proxy).
\end{itemize}

\end{document}
